\documentclass[12pt]{article}
\usepackage[utf8]{inputenc}
\usepackage[T1]{fontenc}
\usepackage{hyperref}
\title{Ke'o Language Resources for K-12 Curriculum Development}
\author{Compiled by Claude Sonnet 5.5 (claude-sonnet-5-5) for Andr\'as Kornai}
\date{}
\begin{document}
\maketitle

\section{Introduction}
Ke'o (also K\'eo; ISO 639-3 xxk; Glottocode keoo1238) is an Austronesian language of central Flores, Indonesia, with about 40,000 speakers. It is highly isolating. Louise Baird's fieldwork in 1998 and 2000 gave the first linguistic description. The search results showed academic and archival materials only. They showed no school materials and no published dictionary.

\section{Primary Educational Resources}
No primers, textbooks or school programmes for Ke'o were found in the search results.

\section{Linguistic and Reference Materials}
\begin{itemize}
\item Louise Baird, \emph{A grammar of K\'eo: an Austronesian language of East Nusantara} (ANU thesis): \url{https://openresearch-repository.anu.edu.au/handle/1885/9446}
\item Baird, Illustration of the IPA, Ke'o (Journal of the International Phonetic Association), which gives the speaker estimate and a phonetic description: \url{https://core-cms.cambridgecore.org/core/services/aop-cambridge-core/content/view/DCA34C5739D6011776B40F1ECC9799BC/S0025100302000178a.pdf/keo.pdf}
\end{itemize}

\section{Religious and Community Resources}
None were found in the search results.

\section{Digital and Audio Resources}
PARADISEC holds Baird's fieldwork collection: recordings and handwritten notebooks with Indonesian and English translations. The notebook items are access-restricted.
\begin{itemize}
\item PARADISEC collection description: \url{https://catalog.paradisec.org.au/repository/LRB2/DV}
\item Notebook item NB1: \url{https://catalog.paradisec.org.au/repository/LRB2/NB1}
\item Notebook item NB2: \url{https://catalog.paradisec.org.au/repository/LRB2/NB2}
\item Notebook item NB4: \url{https://catalog.paradisec.org.au/repository/LRB2/NB4}
\end{itemize}

\section{Curriculum Development Context}
Ke'o is spoken by all members of its community, so a mother-tongue curriculum is feasible in principle. The sources found provide only a grammar and archival data. A curriculum would have to start from Baird's grammar, with community speakers and teachers producing a practical orthography, readers and a dictionary. Indonesian will remain the language of instruction alongside it.
\end{document}